%% ============================================================
%% QCCK-M 第四章 中文版
%% 格式参考：Overleaf 项目 6aa379e8fec9886e0960f9e9（article + ctex 单栏中文稿）
%% 内容来源：ch4/第四章修改2.md（用户中文底稿，逐句原样保留）
%% 处理：去掉「第N句：」写作脚手架；公式排成 equation 环境；
%%       「图 X」引用删除；算法块排成 algorithm 环境
%% 编译：xelatex（ctex 需要 xelatex 或 lualatex）
%% ============================================================

% [已改造] 本文件现为 main.tex 的分章正文（导言区已并入 main.tex），由 main 统一编译。


\section{QCCK-M：面向多变量时间序列预测的量子跨变量耦合网络}

第三章提出 QCCK，将变量间的耦合显式度量为可解释的耦合矩阵 $K$。然而，耦合矩阵只有参与变量表示更新，才能进一步作用于预测过程。
为此，本章将 QCCK 以 QCCK Layer 的形式嵌入编码器层之间，构成完整的预测网络 QCCK-M，使显式耦合能够参与跨变量信息融合。
4.1 给出 QCCK-M 的整体结构与数据流，4.2 给出单个 QCCK Layer 的前向算法与网络训练细节。

\subsection{QCCK-M 总览}

QCCK-M 的整体结构与数据流如下：输入的多变量时序数据经 RevIN 与 Tokenize 预处理，成为每个变量的 $d$ 维表示 token；时序主干编码器逐层更新这些 token，QCCK Layer 则嵌入每个编码器层之后，根据显式度量的耦合进行一次跨变量混合；最后经预测头映射并反归一化，输出各变量的预测 $\hat{Y}$。
RevIN 与 Tokenize 是 QCCK-M 的输入预处理环节。
RevIN 把每个变量沿时间轴减去均值、除以标准差，以缓解样本间的分布漂移.
Tokenize 用一个线性层把每个变量的整段历史观测编码为一个 $d$ 维的变量 token。
变量 token 连同提供时间上下文的时间特征 token 一并作为主干编码器的输入。
主干编码器逐层更新变量 token，每个编码器层在层末输出所有变量的隐状态 $H\in\mathbb{R}^{B\times V\times d}$。
主干编码器层通过双向状态空间扫描实现变量间的信息交互，但该交互过程并不显式输出变量两两之间的耦合关系。
为使第三章得到的显式耦合参与每层变量表示更新，QCCK-M 在每个编码器层之后引入一个 QCCK Layer。
QCCK Layer 以该层的隐状态 $H$ 与对齐频谱 $S$ 为输入，把各变量经 QCCE 编码成量子态，用保真度核度量变量两两之间的耦合，得到耦合矩阵 $K$.随后，$K$ 经去除对角自耦合后直接用于跨变量消息聚合，每个变量依据显式耦合强度接收其他变量的信息，并通过残差方式写回，得到融合跨变量信息的更新隐状态 $H'$。
预测头把每个变量的最终隐状态线性映射为长度等于预测长度的预测序列，并经反归一化还原到原始尺度，得到各变量的预测 $\hat{Y}$。
完成前向预测后，QCCK-M 以时域误差与频域监督项之和作为训练目标，并对整网进行端到端优化，损失函数的具体形式见 4.2。
综上，QCCK-M 通过在主干的逐层表示更新过程中引入 QCCK，使显式度量得到的耦合矩阵 $K$ 成为跨变量信息融合的依据。
因此，$K$ 直接参与变量表示更新，同时以显式矩阵形式保留，可用于检查变量间的耦合关系。单个 QCCK Layer 的前向算法与网络训练细节见 4.2。

\subsection{QCCK-M 的算法}

QCCK-M 在每个主干编码器层之后插入结构相同的 QCCK Layer，各层共享由输入序列得到的对齐频谱 $S$，因此本节先给出单个 QCCK Layer 的前向过程，再说明整网的训练目标与参数设置。

\begin{algorithm}[htbp]
\caption{QCCK Layer 前向传播}
\label{alg:qcck-layer-zh}
\begin{algorithmic}[1]
\REQUIRE 主干隐状态 $H\in\mathbb{R}^{B\times V\times d}$，对齐频谱 $S\in\mathbb{R}^{B\times V\times 2M}$
\ENSURE 更新后的隐状态 $H'$，耦合矩阵 $K$
\STATE 对每个变量 $v$，由 QCCE 构造量子态，经两层 $U_v$ 演化得到 $|\psi_v\rangle$
\STATE 计算保真度核 $K[i,j] = |\langle\psi_i|\psi_j\rangle|^2$
\STATE 去除自耦合：$K \leftarrow K - I$
\STATE 消息投影：$H_q = H \cdot W_q$
\STATE 跨变量聚合：$H_p = K \cdot H_q$
\STATE 残差更新：$H' = \mathrm{LN}(H + \gamma \cdot H_p)$
\RETURN $H'$, $K$
\end{algorithmic}
\end{algorithm}

算法第 1、2 步复用第三章定义的 QCCE 与保真度核。
每个变量 $v$ 由当前层隐状态 $h_v$ 与共享对齐频谱 $S_v$ 构造量子态，并经参数化量子线路 $U_v$ 演化得到 $|\psi_v\rangle$；变量 $i$ 与变量 $j$ 的耦合强度由式~(\ref{eq:qkernel}) 的保真度 $K[i,j] = |\langle\psi_i|\psi_j\rangle|^2$ 给出。
本文固定使用 5 个量子比特，因此每个变量的量子态位于 32 维复 Hilbert 空间。
算法第 3 步用于去除变量自身耦合。
由于保真度核的对角元素恒为 1，即 $K[i,i]=1$，该部分仅表示变量自身相似性，而非跨变量关系。
因此，本文将对角元素置零，仅保留变量之间的耦合信息用于后续跨变量传播。
算法第 4 至 5 步完成消息生成与跨变量融合。
每个变量的隐状态先经可学习投影 $W_q$ 得到消息表示 $H_q$，再由耦合矩阵 $K$ 对不同变量消息进行加权聚合，得到跨变量表示 $H_p$。
由于 $K$ 中每个元素直接对应变量间耦合强度，因此强耦合变量能够贡献更多跨变量信息，而弱耦合变量的影响自然受到抑制。
算法第 6 步通过残差结构完成变量表示更新。
其中，$\gamma$ 为可学习融合系数，用于控制 QCCK Layer 对原始主干表示的注入强度。
通过残差连接，模型能够在保留主干已学习时间模式的同时，引入由量子耦合关系提供的跨变量信息。
量子核需要计算变量两两之间的量子态保真度，其计算复杂度为 $O(V^2\cdot 2^N)$，其中 $V$ 为变量数，$N$ 为量子比特数。
本文固定 $N=5$，因此量子态维度为 32，且该计算不随输入序列长度 $L$ 增长。
QCCK-M 采用时域监督与频域监督联合训练，整体训练目标为：
\begin{equation}
\label{eq:loss-zh}
L = L_{\mathrm{time}} + \lambda \cdot L_{\mathrm{FreDF}},
\end{equation}
其中
\begin{equation}
L_{\mathrm{time}} = \lVert \hat{Y} - Y \rVert^2
\end{equation}
用于约束预测序列与真实序列在时域上的误差；
\begin{equation}
L_{\mathrm{FreDF}} = \mathrm{mean}\,\lvert \mathrm{rFFT}(\hat{Y}) - \mathrm{rFFT}(Y) \rvert
\end{equation}
用于约束二者在频域表示上的差异，$\lambda$ 为频域监督项的权重。
QCCE、参数化量子线路、QCCK Layer 与时序主干中的可学习参数均由上述联合损失进行端到端优化。
QCCK-M 保留主干原有的编码器与预测头结构，仅在编码器层之间引入 QCCK Layer，并增加频域监督项。


